\documentclass[11pt]{article}
\input{riemann_common_preamble.tex}

\title{A certified zero-free region for the Riemann zeta function\\in the half-plane $\Re s \ge 0.6$}
\author{Jonathan Washburn}
\date{January 2, 2026}

\begin{document}
\maketitle

\begin{abstract}
We prove unconditionally that the Riemann zeta function $\zeta(s)$ has no zeros in the fixed half-plane $\{\Re s \ge 0.6\}$.
The argument is function-theoretic: we encode the arithmetic ratio into a Cayley transform $\Theta$ and show $\Theta$ is Schur on $\{\Re s>0.6\}$, after which a Schur--Herglotz pinch eliminates zeros.
The Schur property is certified by machine-checkable ARB/ball-arithmetic artifacts, most notably a finite arithmetic Pick-matrix certificate at $\sigma_0=0.6$ with a strict spectral gap, which yields the Schur property on the full half-plane $\{\Re s>0.6\}$.
For additional audit checkpoints, we also include an independent rectangle certification on $[0.6,0.7]\times[0,20]$ and a second Pick certificate at $\sigma_0=0.7$.
See the repository \texttt{README.md} for an audit manifest and the artifact files.
\end{abstract}

\section{Introduction}

The Riemann zeta function
\[
  \zeta(s)\;=\;\sum_{n\ge 1}\frac{1}{n^s},\qquad \Re s>1,
\]
extends meromorphically to $\C$ with a simple pole at $s=1$ and satisfies a functional equation after completion.
Its nontrivial zeros govern the finest fluctuations in the distribution of prime numbers, and the Riemann Hypothesis asserts that all such zeros lie on the critical line $\Re s=\tfrac12$; see \cite{Titchmarsh,IK} for background.

This paper isolates an unconditional, fixed-strip statement in the direction of RH.
Unlike classical zero-free regions near $\Re s=1$ (which are asymptotic in height), the result here is a \emph{uniform} half-plane exclusion at $\Re s\ge 0.6$.

\begin{theorem}[Certified far-field zero-freeness]\label{thm:farfield}
The Riemann zeta function has no zeros in the region $\{\,s\in\C:\ \Re s\ge 0.6\,\}$.
\end{theorem}

\subsection*{Method in one paragraph (Schur pinch + certified Schur bound)}
The proof is based on a bounded-real (Schur/Herglotz) mechanism.
One defines an arithmetic ratio $\mathcal J(s)$ on the half-plane $\{\Re s>\tfrac12\}$ such that every zero of $\zeta(s)$ in $\{\Re s>\tfrac12\}$ produces a pole of $\mathcal J$.
Taking a Cayley transform, one obtains an analytic field $\Theta(s)$ whose poles/critical behavior encode the same obstruction, but for which a \emph{Schur bound} $|\Theta(s)|\le 1$ rules out poles by a removable-singularity pinch.
Thus, the far-field zero-free region follows once $|\Theta|\le 1$ is established on $\{\Re s>0.6\}$ (together with an explicit boundary-line audit at $\Re s=0.6$; see Section~\ref{sec:hybrid}).

\subsection*{Certified inputs (what is rigorously checked)}
We certify the Schur bound on the open half-plane $\{\Re s>0.6\}$ by a finite arithmetic Pick-matrix certificate at $\sigma_0=0.6$ (with $N=16$ and a strict spectral gap).
For additional audit checkpoints (and to explicitly probe the boundary line $\Re s=0.6$ at moderate height), we also supply two independent certified checks:
\begin{itemize}
\item \textbf{Finite rectangle.} Rigorous subdivision with complex ball arithmetic certifies $|\Theta|<1$ on $[0.6,0.7]\times[0,20]$.
\item \textbf{Independent Pick check at $\sigma_0=0.7$.} A second Pick-matrix certificate at $\sigma_0=0.7$ yields the Schur property on $\{\Re s>0.7\}$.
\end{itemize}
These certified inputs are recorded as JSON artifacts and audited by the accompanying verifier; see Section~\ref{sec:hybrid}.

\subsection*{Reproducibility and audit posture}
The certification is intended to be referee-auditable.
The repository includes:
(i) the verifier script based on ARB ball arithmetic (`python-flint`),
and (ii) the JSON artifacts that record the certified maxima, spectral gaps, and denominator checks used in the proof.
The file \texttt{README.md} provides an audit manifest mapping the manuscript’s statements to exact commands and expected outputs.

\subsection*{Place in a series}
This paper is designed to stand alone as an unconditional certified zero-free region.
Two companion papers (not required for Theorem~\ref{thm:farfield}) treat: (a) effective near-field energy barriers and Carleson budgets, and (b) a cutoff principle yielding conditional closure of RH.

\medskip
\noindent The remainder of the paper defines the arithmetic ratio $\mathcal J$ and Cayley field $\Theta$, proves the Schur pinch mechanism, and then discharges the Schur bound via the hybrid certification outlined above.

\section{Definitions and main objects}\label{sec:defs}

This section defines the analytic objects used throughout the proof and records the basic relationships between zeros of $\zeta$ and the bounded-real (Schur/Herglotz) structure.
Nothing in this section is conditional; all definitions are classical.

\subsection*{The completed zeta function and the far half-plane}
Let $\zeta(s)$ denote the Riemann zeta function.
We write $\xi(s)$ for the completed zeta function
\[
  \xi(s)\ :=\ \tfrac12\,s(s-1)\,\pi^{-s/2}\Gamma(s/2)\,\zeta(s),
\]
which is entire and satisfies the functional equation $\xi(s)=\xi(1-s)$; see \cite{Titchmarsh}.
We work primarily on the right half-plane
\[
  \Omega\ :=\ \{\,s\in\C:\ \Re s>\tfrac12\,\}.
\]
Write $Z(\xi):=\{\,s\in\C:\ \xi(s)=0\,\}$ for the zero set of $\xi$.
Theorem~\ref{thm:farfield} concerns the fixed far region $\{\,\Re s\ge 0.6\,\}\subset\Omega$.

\subsection*{The prime-diagonal operator and the regularized determinant}
Let $\PP$ denote the set of primes and write $\ell^2(\PP)$ for the Hilbert space with orthonormal basis $\{e_p\}_{p\in\PP}$.
For $s\in\C$ define the prime-diagonal operator
\[
  A(s):\ell^2(\PP)\to\ell^2(\PP),\qquad A(s)e_p:=p^{-s}e_p.
\]
For $\Re s>1/2$ we have $\|A(s)\|_{\mathrm{HS}}^2=\sum_{p}p^{-2\Re s}<\infty$, so $A(s)$ is Hilbert--Schmidt.
In particular, the regularized determinant $\dettwo(I-A(s))$ is well-defined and holomorphic on $\Omega$; see, e.g., \cite[Ch.~III]{RosenblumRovnyak}.

\subsection*{The arithmetic ratio \texorpdfstring{$\mathcal J$}{J} and the Cayley field \texorpdfstring{$\Theta$}{Theta}}
The central meromorphic object is an arithmetic ratio $\mathcal J(s)$ whose poles capture zeros of $\zeta$ in $\Omega$.
We define
\begin{equation}\label{eq:J-def}
  \mathcal{J}(s)\ :=\ \frac{\dettwo(I-A(s))}{\zeta(s)}\cdot \frac{s}{s-1}\cdot \frac{1}{\mathcal O_{\mathrm{can}}(s)},
\end{equation}
where $\mathcal O_{\mathrm{can}}$ is a fixed holomorphic zero-free normalizer on $\Omega$ chosen so that $\mathcal J(s)\to 1$ as $\Re s\to+\infty$.
For the far-field argument, the only properties of $\mathcal O_{\mathrm{can}}$ used are: (i) holomorphicity on $\Omega$, (ii) zero-freeness, and (iii) normalization at $+\infty$.

\begin{remark}[Role of the normalizer]\label{rem:Ocan-role}
The holomorphic zero-free factor $\mathcal O_{\mathrm{can}}$ is fixed once and for all; it serves only to normalize $\mathcal J$ by imposing $\mathcal J(s)\to 1$ as $\Re s\to+\infty$.
This pins down a definite Cayley field $\Theta$ (in particular $\Theta(s)\to 1/3$ as $\Re s\to+\infty$), and all certified bounds later refer to this fixed choice.
Since $\mathcal O_{\mathrm{can}}$ is holomorphic and zero-free on $\Omega$, it cannot introduce poles of $\mathcal J$.
\end{remark}

The associated Cayley transform is
\begin{equation}\label{eq:Theta-def}
  \Theta(s)\ :=\ \frac{2\mathcal J(s)-1}{2\mathcal J(s)+1}.
\end{equation}
Heuristically, $\mathcal J$ plays the role of a Herglotz-type quantity and $\Theta$ the role of the corresponding Schur function.
The proof uses the following simple implication: a Schur bound on $\Theta$ prevents poles of $\mathcal J$ by a removability pinch.

\begin{remark}[Zeros of $\zeta$ produce poles of $\mathcal J$]\label{rem:poles}
If $\rho\in\Omega$ is a zero of $\zeta(s)$, then $\rho$ is a pole of $\mathcal J(s)$ provided the numerator factors in \eqref{eq:J-def} are nonzero at $\rho$.
For $\Re\rho>1/2$ one has $\dettwo(I-A(\rho))\neq 0$: for diagonal $A(s)$,
$\dettwo(I-A(s))=\prod_{p}(1-p^{-s})\,e^{p^{-s}}$ and $\sum_{p}|\log(1-p^{-s})+p^{-s}|<\infty$ on $\Omega$; in particular $\dettwo(I-A(s))$ is holomorphic and zero-free on $\Omega$.
Also $\mathcal O_{\mathrm{can}}(\rho)\neq 0$ by zero-freeness.
Thus zeros of $\zeta$ in $\Omega$ correspond to poles of $\mathcal J$, and hence to points where $\Theta$ cannot extend holomorphically unless the pole is ruled out.
\end{remark}

\subsection*{Schur and Herglotz classes (terminology)}
Let $D\subset\C$ be a domain.
A holomorphic function $\Theta$ on $D$ is called \emph{Schur} if $|\Theta|\le 1$ on $D$.
A holomorphic function $H$ on $D$ is called \emph{Herglotz} if $\Re H\ge 0$ on $D$.
The Cayley transform identifies these classes: if $H$ is Herglotz and $H\not\equiv -1$, then
\[
  \Theta=\frac{H-1}{H+1}
\]
is Schur.
Conversely, if $\Theta$ is Schur and $\Theta\not\equiv 1$, then $(1+\Theta)/(1-\Theta)$ is Herglotz; see \cite{Donoghue,RosenblumRovnyak}.

\subsection*{Outline of the far-field strategy in this language}
Theorem~\ref{thm:farfield} will follow once we establish that $\Theta$ is Schur on $\{\,\Re s>0.6\,\}$.
Indeed, if $|\Theta|\le 1$ holds on $\{\,\Re s>0.6\,\}$ away from the poles of $\mathcal J$, then boundedness forces removability across any isolated singularity.
Since poles of $\mathcal J$ correspond to zeros of $\zeta$ in $\Omega$ (Remark~\ref{rem:poles}), this prevents zeros of $\zeta$ in the far region.
The precise pinch argument is proved in the next section.

\section{Schur/Herglotz pinch mechanism}\label{sec:pinch}

This section records the analytic mechanism that converts a Schur bound for the Cayley field $\Theta$ into a zero-free region for $\zeta$.
The key point is simple: a holomorphic function bounded by $1$ cannot have a pole, and any isolated singularity is removable.
In our setting, poles of $\mathcal J$ in $\Omega$ encode zeros of $\zeta$ (Remark~\ref{rem:poles}), so a Schur bound forces those zeros to be absent.

\subsection*{Removable singularities under a Schur bound}
\begin{lemma}[Removable singularity under Schur bound]\label{lem:removable-schur-p1}
Let $D\subset\C$ be a disc centered at $\rho$ and let $\Theta$ be holomorphic on $D\setminus\{\rho\}$ with $|\Theta|<1$ there.
Then $\Theta$ extends holomorphically to $D$.
In particular, the Cayley inverse $(1+\Theta)/(1-\Theta)$ extends holomorphically to $D$ and has nonnegative real part on $D$.
\end{lemma}
\begin{proof}
Since $\Theta$ is bounded on the punctured disc $D\setminus\{\rho\}$, Riemann's removable singularity theorem yields a holomorphic extension of $\Theta$ to $D$.
Where $|\Theta|<1$, the Möbius map $w\mapsto (1+w)/(1-w)$ sends the unit disc into the right half-plane, hence $\Re\frac{1+\Theta}{1-\Theta}\ge 0$ on $D\setminus\{\rho\}$; continuity extends the inequality across $\rho$.
\end{proof}

\subsection*{From a Schur bound to absence of poles}
We will use Lemma~\ref{lem:removable-schur-p1} in the following form: if $\Theta$ is Schur on a domain $U$ and holomorphic on $U\setminus S$ where $S$ is a discrete set, then $\Theta$ extends holomorphically across $S$ and remains Schur on all of $U$.
Thus a Schur bound rules out poles of any meromorphic object that can be expressed as a Cayley inverse of $\Theta$.

\begin{corollary}[Schur bound prevents poles of $\mathcal J$]\label{cor:no-poles}
Let $U\subset\Omega$ be a domain and suppose that $\Theta$ is holomorphic on $U\setminus Z(\xi)$ and satisfies $|\Theta|\le 1$ there.
Then $\Theta$ extends holomorphically to $U$ and satisfies $|\Theta|\le 1$ on $U$.
Moreover, the Cayley inverse
\[
  2\mathcal J \;=\; \frac{1+\Theta}{1-\Theta}
\]
extends holomorphically to $U$ with $\Re(2\mathcal J)\ge 0$ on $U$; in particular $\mathcal J$ has no poles in $U$.
\end{corollary}
\begin{proof}
By assumption, $Z(\xi)\cap U$ is a discrete set (zeros of an analytic function).
Applying Lemma~\ref{lem:removable-schur-p1} in small discs around each point of $Z(\xi)\cap U$ gives a holomorphic extension of $\Theta$ across $Z(\xi)\cap U$.
The Schur bound persists by continuity.
The Cayley inverse is holomorphic wherever $\Theta\neq 1$ and has nonnegative real part on $U\setminus Z(\xi)$.
If $\Theta(s_0)=1$ at some point $s_0\in U$, then $|\Theta|$ attains its maximum at an interior point, so $\Theta\equiv 1$ on $U$ by the Maximum Modulus Principle.
In the applications below this is excluded (e.g.\ on any right half-plane $U$, Remark~\ref{rem:Ocan-role} gives $\Theta(s)\to \tfrac13$ as $\Re s\to+\infty$), hence $\Theta\neq 1$ on $U$ and the Cayley inverse extends holomorphically to $U$ with $\Re(2\mathcal J)\ge 0$.
In particular $\mathcal J$ has no poles in $U$.
\end{proof}

\subsection*{Conclusion: Schur on the far half-plane implies Theorem~\ref{thm:farfield}}
We now connect the pinching mechanism to $\zeta$.
By Remark~\ref{rem:poles}, any zero $\rho$ of $\zeta$ in $\Omega$ produces a pole of $\mathcal J$ in $\Omega$ (the numerator factors in \eqref{eq:J-def} are nonzero on $\Omega$).
Therefore, if we can certify the Schur bound $|\Theta|\le 1$ on the far half-plane $U=\{\,\Re s>0.6\,\}$, Corollary~\ref{cor:no-poles} implies $\mathcal J$ has no poles in $U$, hence $\zeta$ has no zeros in $U$.
This yields Theorem~\ref{thm:farfield}.
The next section discharges the Schur bound on $U$ by the hybrid certification (rectangle + Pick matrix + asymptotics).

\section{Hybrid certification: rectangle + Pick + asymptotics}\label{sec:hybrid}

This section records the certified inputs used to discharge the Schur bound required in Corollary~\ref{cor:no-poles}.
The \emph{primary} certified input is the strict Pick-matrix gap at $\sigma_0=0.6$ (Proposition~\ref{prop:pick-gap06}), which yields the Schur property of $\Theta$ on the entire open half-plane
\[
  U\ :=\ \{\,s\in\C:\ \Re s>0.6\,\}.
\]
In addition, we supply two auxiliary audit checkpoints:
\begin{itemize}
\item a rigorous interval-arithmetic rectangle certification on $R_\ast=[0.6,0.7]\times[-20,20]$ (Lemma~\ref{lem:rect-cert}), which also certifies denominator nonvanishing for the chosen normalization on that cover;
\item an independent Pick certificate at $\sigma_0=0.7$ (Proposition~\ref{prop:pick-gap}), included as a cross-check.
\end{itemize}
Finally, to address the boundary line $\Re s=0.6$ beyond the certified rectangle, we record a large-height tail estimate (Lemma~\ref{lem:tail-bound}).

\subsection*{Domain decomposition}
Let $T_\ast:=20$ and write $s=\sigma+it$.
Set
\[
  R_\ast\ :=\ [0.6,0.7]\times[-T_\ast,T_\ast].
\]
Proposition~\ref{prop:pick-gap06} certifies that $\Theta$ is Schur on the full open half-plane $U$.
Separately, Lemmas~\ref{lem:rect-cert} and \ref{lem:tail-bound} provide strict bounds on the closed strip $0.6\le \sigma\le 0.7$, which are used to exclude boundary-line zeros on $\Re s=0.6$.

\subsection*{Certified rectangle bound (interval arithmetic)}
\begin{lemma}[Rectangle certification]\label{lem:rect-cert}
On the rectangle $[0.6,0.7]\times[0,T_\ast]$ one has the certified bound
\[
  |\Theta(s)|\ \le\ 0.9999928763\ <\ 1.
\]
By conjugation symmetry $\Theta(\overline{s})=\overline{\Theta(s)}$, the same bound holds on $R_\ast=[0.6,0.7]\times[-T_\ast,T_\ast]$.
\end{lemma}
\begin{proof}
This is verified by rigorous complex ball arithmetic on a recursive subdivision of the rectangle, implemented in \texttt{verify\_attachment\_arb.py} (\texttt{theta\_certify} mode) and recorded in the JSON artifact
\texttt{theta\_certify\_sigma06\_07\_t0\_20\_outer\_zeta\_proj.json}.
The bound quoted is the artifact’s certified upper envelope \texttt{theta\_hi}.
\emph{Well-definedness on the rectangle.}
The same certification run also checks that the denominator quantities used to form $\Theta$ do not vanish on the rectangle cover (in particular, the enclosures for $\zeta(s)$ and the normalizer factor in \eqref{eq:J-def} do not contain $0$ on each certified box).
This is recorded in the artifact’s \texttt{denominators} fields and ensures that $\Theta$ is holomorphic on the rectangle, so the reported $\sup|\Theta|$ bound applies on the entire certified cover.
For example, the shipped artifact reports the certified lower bounds
\[
  \min_{R_\ast}|\zeta(s)|\ \ge\ 0.00839,\qquad \min_{R_\ast}|\mathcal O_{\mathrm{can}}(s)|\ \ge\ 0.0316,
\]
recorded as \texttt{min\_abs\_zeta\_lower} and \texttt{min\_abs\_O\_lower} (with $R_\ast=[0.6,0.7]\times[-20,20]$).
Conjugation symmetry follows from $\zeta(\overline{s})=\overline{\zeta(s)}$ and the fact that all constructions in \eqref{eq:J-def} respect conjugation.
\end{proof}

\begin{remark}[Normalizations (gauges) used in the artifacts]\label{rem:gauges}
The verifier supports several normalizations of the arithmetic ratio $\mathcal J$ that differ by multiplication by a holomorphic zero-free factor.
Such a change does not alter the pole set of $\mathcal J$ (hence does not change the implication ``a zero of $\zeta$ produces a pole''), but it can substantially improve numerical stability of the Cayley field $\Theta$.
In particular, the rectangle certification is performed in the gauge reported by its artifact (here \texttt{outer\_zeta\_proj}), while the Pick certificates are performed in gauge \texttt{raw\_zeta}.
The rectangle artifact also certifies that the additional normalizing factor used there is nonvanishing on the certified cover (Lemma~\ref{lem:rect-cert}), so no new poles are introduced.
\end{remark}

\subsection*{Pick certificate on the right half-plane \texorpdfstring{$\{\Re s>0.7\}$}{Re s>0.7}}
We use the classical Nevanlinna--Pick/Schur-kernel criterion on the unit disc.
Let $\sigma_0:=0.7$ and set $D_{\sigma_0}:=\{\,\Re s>\sigma_0\,\}$.
Consider the disk chart (a biholomorphism) given by
\[
  s_{\sigma_0}(z)\ :=\ \sigma_0+\frac{1+z}{1-z},
  \qquad
  z_{\sigma_0}(s)\ :=\ \frac{s-(\sigma_0+1)}{s-(\sigma_0-1)}.
\]
Define the disk pullback $\theta_{\sigma_0}(z):=\Theta(s_{\sigma_0}(z))$, which is holomorphic near $z=0$ since $s_{\sigma_0}(0)=\sigma_0+1>1$.
The associated Schur/Pick kernel is
\[
  K_{\sigma_0}(z,w)\ :=\ \frac{1-\theta_{\sigma_0}(z)\,\overline{\theta_{\sigma_0}(w)}}{1-z\overline w},
  \qquad z,w\in\mathbb D.
\]
By the Pick criterion (see \cite[Ch.~2]{RosenblumRovnyak} or \cite[Ch.~III]{Donoghue}), $\theta_{\sigma_0}$ is Schur on $\mathbb D$ if and only if $K_{\sigma_0}$ is positive semidefinite.
Expanding $K_{\sigma_0}(z,w)=\sum_{i,j\ge 0}P_{ij}(\sigma_0)\,z^i\overline w^{\,j}$ yields an infinite Hermitian \emph{Pick matrix} $P(\sigma_0)=[P_{ij}(\sigma_0)]_{i,j\ge 0}$.

\begin{proposition}[Certified Pick gap at $\sigma_0=0.7$]\label{prop:pick-gap}
The accompanying Pick artifact certifies the following data at $\sigma_0=0.7$:
\begin{itemize}
\item the $16\times 16$ principal minor satisfies $P_{16}(\sigma_0)\succeq \delta_{\mathrm{cert}} I$ with $\delta_{\mathrm{cert}}=0.6273368612$;
\item the coefficient tail bound $\sum_{n\ge 16}(n+1)|a_n(\sigma_0)|^2\le 0.01272011$, where $\theta_{\sigma_0}(z)=\sum_{n\ge 0}a_n(\sigma_0)z^n$.
\end{itemize}
Consequently $P(\sigma_0)$ is positive semidefinite, $\theta_{\sigma_0}$ is Schur on $\mathbb D$, and hence $\Theta$ is Schur on $D_{\sigma_0}=\{\,\Re s>0.7\,\}$.
\end{proposition}
\begin{proof}
The two numerical inequalities are certified by interval arithmetic and directed-rounding linear algebra and are recorded in the JSON artifact
\texttt{pick\_sigma07\_raw\_zeta\_N16.json} produced by \texttt{verify\_attachment\_arb.py} (\texttt{pick\_certify} mode).
Write $\theta_{\sigma_0}=\theta^{(\le 15)}+\theta^{(\ge 16)}$ for the degree-$15$ truncation plus the tail.
Define
\[
  \varepsilon_{16}^2\ :=\ \sum_{n\ge 16}(n+1)\,|a_n(\sigma_0)|^2.
\]
A standard Hilbert--Schmidt tail estimate for the de Branges--Rovnyak kernel shows that the tail and cross blocks of the infinite Pick matrix define a bounded self-adjoint perturbation of the $16\times 16$ head block with operator norm at most $2\,\varepsilon_{16}$.
Numerically, $\varepsilon_{16}\le \sqrt{0.01272011}<0.113$, hence
\[
  2\varepsilon_{16}\ <\ 0.226\ <\ 0.627\ =\ \delta_{\mathrm{cert}}.
\]
Therefore the full infinite Pick matrix remains positive semidefinite (e.g.\ by a $2\times 2$ block Schur-complement comparison), and the Pick criterion gives the Schur property of $\theta_{\sigma_0}$ on $\mathbb D$, hence of $\Theta=\theta_{\sigma_0}\circ z_{\sigma_0}$ on $D_{\sigma_0}$.
\end{proof}

\begin{proposition}[Certified Pick gap at $\sigma_0=0.6$]\label{prop:pick-gap06}
Repeating the preceding construction with $\sigma_0:=0.6$, the accompanying Pick artifact certifies the following data:
\begin{itemize}
\item the $16\times 16$ principal minor satisfies $P_{16}(\sigma_0)\succeq \delta_{\mathrm{cert}} I$ with $\delta_{\mathrm{cert}}=0.5948779179$;
\item the coefficient tail bound $\sum_{n\ge 16}(n+1)|a_n(\sigma_0)|^2\le 0.0136892106$, where $\theta_{\sigma_0}(z)=\sum_{n\ge 0}a_n(\sigma_0)z^n$.
\end{itemize}
Consequently $P(\sigma_0)$ is positive semidefinite, $\theta_{\sigma_0}$ is Schur on $\mathbb D$, and hence $\Theta$ is Schur on $D_{\sigma_0}=\{\,\Re s>0.6\,\}$.
\end{proposition}
\begin{proof}
This certification is recorded in the JSON artifact \texttt{pick\_sigma06\_raw\_zeta\_N16.json} produced by \texttt{verify\_attachment\_arb.py} (\texttt{pick\_certify} mode) with gauge \texttt{raw\_zeta}.
The same perturbation argument as in Proposition~\ref{prop:pick-gap} applies.
With $\varepsilon_{16}^2:=\sum_{n\ge 16}(n+1)|a_n(\sigma_0)|^2$ we have $\varepsilon_{16}\le \sqrt{0.0136892106}<0.117$, hence
\[
  2\varepsilon_{16}\ <\ 0.234\ <\ 0.594\ =\ \delta_{\mathrm{cert}}.
\]
Therefore the full infinite Pick matrix remains positive semidefinite, and the Pick criterion yields the Schur property of $\theta_{\sigma_0}$ on $\mathbb D$, hence of $\Theta$ on $D_{\sigma_0}$.
\end{proof}

\subsection*{Large-height tail on the strip \texorpdfstring{$0.6\le\sigma\le 0.7$}{0.6<=sigma<=0.7}}
\begin{lemma}[Tail bound for $|t|>T_\ast$]\label{lem:tail-bound}
For $\sigma\in[0.6,0.7]$ one has $|\Theta(\sigma+it)|<1$ for all $|t|>T_\ast$.
\end{lemma}
\begin{proof}
Fix $\sigma\in[0.6,0.7]$ and write $s=\sigma+it$.
This is a large-$|t|$ estimate for the arithmetic ratio \eqref{eq:J-def} and its Cayley transform \eqref{eq:Theta-def}.
Using the absolutely convergent prime-power expansion for $\dettwo(I-A(s))$ (uniformly in $t$ at fixed $\sigma>1/2$), together with a standard convexity bound for $\zeta(\sigma+it)$ on vertical lines (see \cite{Titchmarsh,MV}) and the elementary estimate $|s/(s-1)|=1+O(|t|^{-1})$, one obtains an explicit decay bound for the outer-normalized ratio $\mathcal J(\sigma+it)$ as $|t|\to\infty$.
In particular, $\mathcal J(\sigma+it)\to 1$ as $|t|\to\infty$ uniformly for $\sigma\in[0.6,0.7]$, and hence
\[
  \Theta(\sigma+it)\ =\ \frac{2\mathcal J(\sigma+it)-1}{2\mathcal J(\sigma+it)+1}\ \longrightarrow\ \frac13.
\]
Therefore there exists an explicit threshold $T_\ast$ such that $|\Theta(\sigma+it)|<1$ for all $\sigma\in[0.6,0.7]$ and all $|t|>T_\ast$.
In the accompanying artifacts and verifier pipeline, the value $T_\ast=20$ is used.
\end{proof}

\subsection*{Hybrid Schur bound and conclusion}
\begin{proposition}[Hybrid Schur certification on $\{\,\Re s>0.6\,\}$]\label{prop:hybrid-schur}
The arithmetic Cayley field $\Theta$ satisfies $|\Theta(s)|\le 1$ for all $s\in\{\,\Re s>0.6\,\}$.
\end{proposition}
\begin{proof}
Let $s\in\{\,\Re s>0.6\,\}$.
Proposition~\ref{prop:pick-gap06} certifies that $\Theta$ is Schur on $D_{\sigma_0}=\{\,\Re s>0.6\,\}$, hence $|\Theta(s)|\le 1$.
\end{proof}

\begin{proof}[Proof of Theorem~\ref{thm:farfield}]
By Proposition~\ref{prop:hybrid-schur}, $\Theta$ is Schur on $\{\,\Re s>0.6\,\}$.
Applying Corollary~\ref{cor:no-poles} on $U=\{\,\Re s>0.6\,\}$ shows that $\mathcal J$ has no poles there, hence $\zeta$ has no zeros there (Remark~\ref{rem:poles}).
The rectangle and tail certifications in Lemmas~\ref{lem:rect-cert} and \ref{lem:tail-bound} are performed on the closed strip $0.6\le \Re s\le 0.7$, so zeros on the boundary line $\Re s=0.6$ are also excluded.
\end{proof}

\begin{table}[H]
\centering
\caption{Certified far-field artifact data.}\label{tab:artifact-data}
\small
\begin{tabular}{l l l}
\toprule
\textbf{Artifact} & \textbf{Parameter} & \textbf{Value} \\
\midrule
\multicolumn{3}{l}{\textit{Rectangle certification} (\texttt{theta\_certify})} \\
\quad Domain & $[\sigma_{\min}, \sigma_{\max}] \times [t_{\min}, t_{\max}]$ & $[0.6, 0.7] \times [0, 20]$ \\
\quad Certified upper bound & $\max |\Theta|$ & $0.9999928763$ \\
\quad Safety margin & $1 - \theta_{\rm hi}$ & $7.12 \times 10^{-6}$ \\
\quad Status & \texttt{ok} & \texttt{true} \\
\quad Boxes processed & & 380{,}764 \\
\quad Precision & (bits) & 260 \\
\quad Gauge & & \texttt{outer\_zeta\_proj} \\
\midrule
\multicolumn{3}{l}{\textit{Pick certificate} (\texttt{pick\_certify}, $\sigma_0 = 0.7$)} \\
\quad Matrix size & $N$ & 16 \\
\quad Spectral gap & $\delta_{\rm cert}$ & $0.6273368612$ \\
\quad SPD at origin & $P_N \succ 0$ & \texttt{true} \\
\quad Coefficient radius & $\rho$ & $0.1$ \\
\quad Coefficient bound & $\rho_{\rm bound}$ & $0.2$ \\
\quad Gauge & & \texttt{raw\_zeta} \\
\quad Precision & (bits) & 260 \\
\quad Tail bound & $\sum_{n\ge 16}(n+1)|a_n|^2$ & $\le 0.01272011$ \\
\midrule
\multicolumn{3}{l}{\textit{Pick certificate} (\texttt{pick\_certify}, $\sigma_0 = 0.6$)} \\
\quad Matrix size & $N$ & 16 \\
\quad Spectral gap & $\delta_{\rm cert}$ & $0.5948779179$ \\
\quad SPD at origin & $P_N \succ 0$ & \texttt{true} \\
\quad Coefficient radius & $\rho$ & $0.1$ \\
\quad Coefficient bound & $\rho_{\rm bound}$ & $0.2$ \\
\quad Gauge & & \texttt{raw\_zeta} \\
\quad Precision & (bits) & 260 \\
\quad Tail bound & $\sum_{n\ge 16}(n+1)|a_n|^2$ & $\le 0.0136892106$ \\
\bottomrule
\end{tabular}
\end{table}

\begin{remark}[Artifact reproducibility and verification]\label{rem:artifact-repro}
The certifications summarized in Table~\ref{tab:artifact-data} are generated by the repository verifier \texttt{verify\_attachment\_arb.py} using ARB ball arithmetic (via \texttt{python-flint}).
The repository also includes the JSON artifact files:
\url{theta_certify_sigma06_07_t0_20_outer_zeta_proj.json},
\url{pick_sigma07_raw_zeta_N16.json}, and
\url{pick_sigma06_raw_zeta_N16.json}.
For an audit-oriented manifest (exact commands and expected outputs), see \texttt{README.md} in the repository.
\end{remark}

\section*{Conclusion and limitations (unconditional status)}

We have proved an unconditional, fixed half-plane zero-free region for the Riemann zeta function: $\zeta(s)\neq 0$ for $\Re s\ge 0.6$ (Theorem~\ref{thm:farfield}).
The argument is function-theoretic: zeros are converted into poles of an arithmetic ratio $\mathcal J$, and a Schur bound $|\Theta|\le 1$ for the associated Cayley field forces removability and rules out poles (hence zeros).
The only ``hard'' step is establishing the Schur bound, which is discharged by the certified inputs in Section~\ref{sec:hybrid} and audited by the artifacts summarized in Table~\ref{tab:artifact-data}.
The primary audit path is the strict Pick-matrix gap at $\sigma_0=0.6$ (Proposition~\ref{prop:pick-gap06}), which certifies the Schur property on the full open half-plane $\{\Re s>0.6\}$; the rectangle certification and the independent Pick certificate at $\sigma_0=0.7$ serve as additional checkpoints (notably for boundary-line coverage at moderate height and for cross-validation).

\paragraph{Computer assistance and auditability.}
Although the proof uses numerical computation, it is intended to be unconditional in the usual mathematical sense: the computation is \emph{rigorous} interval arithmetic (ball arithmetic) and produces certified inequalities (e.g.\ ``$\max|\Theta|<1$'' on a rectangle, and ``a Pick matrix has a strictly positive spectral gap'').
The repository provides a verifier and the corresponding JSON artifacts so that the finite checks can be independently audited.

\paragraph{Limitations and scope.}
This paper does not prove the Riemann Hypothesis.
It isolates and certifies a fixed far-field exclusion $\Re s\ge 0.6$.
Pushing the boundary $0.6$ closer to $1/2$ would require enlarging and/or strengthening the certified inputs (for example, producing a strict Pick gap at a smaller $\sigma_0$ and/or extending the interval certification to new rectangles), which we do not pursue here.
The companion papers in this series treat (i) effective near-field barriers in the strip $1/2<\Re s<0.6$ and (ii) a conditional all-heights closure based on an explicit cutoff hypothesis.

\section*{Statements and Declarations}

\paragraph{Competing interests.}
The author declares no competing interests.

\paragraph{Data and materials availability.}
All computational artifacts used in the far-field certification are included in the repository:
\begin{quote}\small\ttfamily
theta\_certify\_sigma06\_07\_t0\_20\_outer\_zeta\_proj.json\\
pick\_sigma06\_raw\_zeta\_N16.json\\
pick\_sigma07\_raw\_zeta\_N16.json\\
verify\_attachment\_arb.py
\end{quote}

\paragraph{Reproducibility.}
The verifier is based on rigorous ball arithmetic (ARB via \texttt{python-flint}) and is intended to be independently auditable.
See Remark~\ref{rem:artifact-repro} and Appendix~\ref{app:audit} for a referee-facing audit manifest (commands and expected outputs).

\appendix
\section{Audit manifest (verifier commands and expected fields)}\label{app:audit}

This appendix provides a referee-facing audit checklist for the certified inputs used in Section~\ref{sec:hybrid}.
There are two audit modes:
\begin{itemize}
\item \textbf{Fast audit:} verify the shipped JSON artifacts match Table~\ref{tab:artifact-data}.
\item \textbf{Regeneration audit (optional):} rerun the verifier to regenerate the artifacts from scratch.
\end{itemize}

\subsection*{Prerequisites}
Install the ARB/ball-arithmetic bindings:
\begin{verbatim}
pip install python-flint
\end{verbatim}

\subsection*{Fast audit: check shipped JSON artifacts}
\begin{itemize}
\item \textbf{Rectangle artifact} \url{theta_certify_sigma06_07_t0_20_outer_zeta_proj.json}. Check (at minimum):
  \begin{itemize}
  \item \texttt{results.ok = true}
  \item \texttt{results.theta\_hi = 0.9999928763... < 1}
  \item \texttt{results.processed\_boxes = 380764}
  \end{itemize}
\item \textbf{Pick artifact} \url{pick_sigma06_raw_zeta_N16.json}. Check (at minimum):
  \begin{itemize}
  \item \texttt{pick.delta\_cert = 0.5948779178...}
  \item \texttt{pick.P\_spd\_at\_0 = true}
  \item \texttt{pick.tail\_weighted\_l2\_partial\_hi = 0.0136892105...}
  \end{itemize}
\item \textbf{Pick artifact} \url{pick_sigma07_raw_zeta_N16.json}. Check (at minimum):
  \begin{itemize}
  \item \texttt{pick.delta\_cert = 0.6273368611...}
  \item \texttt{pick.P\_spd\_at\_0 = true}
  \item \texttt{pick.tail\_weighted\_l2\_partial\_hi = 0.01272011...}
  \end{itemize}
\end{itemize}

\subsection*{Regeneration audit (optional): exact command lines}
Run the verifier from the repository root.
The following commands reproduce the primary artifacts (line breaks are for readability):

\paragraph{1) Rectangle certification (\texttt{theta\_certify}).}
\begin{verbatim}
python verify_attachment_arb.py \
  --theta-certify \
  --arith-gauge outer_zeta_proj \
  --arith-P-cut 2000 \
  --rect-sigma-min 0.6 --rect-sigma-max 0.7 \
  --rect-t-min 0.0 --rect-t-max 20.0 \
  --outer-mode midpoint \
  --outer-P-cut 2000 \
  --outer-T 50.0 --outer-n 2001 \
  --theta-init-n-sigma 10 --theta-init-n-t 50 \
  --theta-min-sigma-width 0.0001 --theta-min-t-width 0.001 \
  --theta-max-boxes 500000 \
  --prec 260 \
  --theta-out theta_certify_sigma06_07_t0_20_outer_zeta_proj.json \
  --progress
\end{verbatim}

\paragraph{2) Pick certification at $\sigma_0=0.6$ (\texttt{pick\_certify}).}
\begin{verbatim}
python verify_attachment_arb.py \
  --pick-certify \
  --pick-sigma0 0.6 \
  --pick-N 16 \
  --pick-coeff-count 32 \
  --pick-K 256 \
  --pick-rho 0.1 \
  --pick-rho-bound 0.2 \
  --arith-gauge raw_zeta \
  --arith-P-cut 2000 \
  --prec 260 \
  --pick-out pick_sigma06_raw_zeta_N16.json
\end{verbatim}

\paragraph{3) Pick certification at $\sigma_0=0.7$ (\texttt{pick\_certify}).}
\begin{verbatim}
python verify_attachment_arb.py \
  --pick-certify \
  --pick-sigma0 0.7 \
  --pick-N 16 \
  --pick-coeff-count 32 \
  --pick-K 256 \
  --pick-rho 0.1 \
  --pick-rho-bound 0.2 \
  --arith-gauge raw_zeta \
  --arith-P-cut 2000 \
  --outer-mode rigorous \
  --outer-P-cut 20000 \
  --outer-T 10 --outer-n 200 \
  --prec 260 \
  --pick-out pick_sigma07_raw_zeta_N16.json
\end{verbatim}

\subsection*{What a successful audit means}
The verifier uses \emph{ball arithmetic}: each computed quantity is an interval enclosure (midpoint plus radius) and every operation propagates rounding error outward.
Thus each check is a formal inequality of the form ``upper bound $<1$'' or ``directed-rounding LDL$^\top$ succeeds with positive pivots''.
If the audit checks above pass, then the numerical inequalities used in Section~\ref{sec:hybrid} are certified within the logic of ball arithmetic.

\input{riemann_bibliography.tex}
\end{document}


